\section*{Ethics Statement}

We affirm compliance with the ICLR Code of Ethics. This work studies self-improving AI systems in the limited context of code-editing agents evaluated on standard programming benchmarks. No human subjects were involved and no personally identifiable information (PII) was collected or processed; IRB approval was therefore not required.

\textbf{Safety and misuse.}
Self-modifying systems can pose safety risks if allowed to act without constraints or if optimizations inadvertently introduce unsafe behaviors. To mitigate this, all agents in our experiments ran inside isolated sandboxes with strict resource and time limits; agents had limited network access and no ability to modify the host environment. The self-improvement scope was restricted to the agent’s own Python codebase and evaluation harnesses. We maintained a complete, auditable lineage (archive) of code changes and evaluations, enabling rollback and post-hoc analysis. We did not deploy discovered agents in real development environments. Our release plan (code, prompts, and evaluation artifacts) will exclude any components that grant elevated system access and will include default sandboxing, guardrails, and clear documentation of intended use.

\textbf{Dual-use, downstream impact, and limitations.}
Stronger autonomous coding agents could be dual-use (e.g., aiding software maintenance, but also potentially facilitating creation of harmful code if misapplied). We believe the research benefits (e.g., advancing methods for controlled, auditable self-improvement and demonstrating practical safeguards) outweigh the risks. Nevertheless, we explicitly discourage security-sensitive or unsandboxed deployment and provide concrete safety recommendations (\Cref{sec:safety}). Our empirical focus on benchmark optimization may not capture all desirable properties (robustness, interpretability, or broader social values). We therefore treat benchmark gains as necessary but insufficient indicators of general AI development, and discuss avenues to integrate other objectives (e.g., safety, reasoning) into the optimization loop.

\textbf{Data governance, IP, and licensing.}
We evaluate on SWE-bench Verified and Polyglot, which are composed of open-source repositories and tasks. We complied with dataset licenses and usage terms to the best of our knowledge. We did not introduce or distribute proprietary code. Foundation models (FMs) were accessed via provider APIs under their terms of service; we did not submit sensitive data, nor attempt to circumvent usage policies. Logs released with this work will be scrubbed of API keys and any incidental sensitive strings.

\textbf{Bias, fairness, and equity.}
Although our domain is software code rather than human-centered text, FM behavior can still reflect biases (e.g., language or ecosystem preferences) and may unevenly benefit communities whose tooling is better represented in training data. We partially address this by evaluating across multiple languages (Polyglot) and reporting cross-benchmark transfer. Future work should add diagnostics for biased failure modes and include broader, community-driven task sets.

\textbf{Conflicts of interest and funding.}
No author has a financial interest in products whose performance is evaluated here. Sponsors and employers did not influence experimental design, analysis, or the decision to publish, beyond providing salary or standard research support. Any external compute or API credits are acknowledged in the appendix.
